%------------------------------------------------------
% Harvard-style one-page resume -- NGUYEN KHAC BAO
% Target: Asoft - AI Automation / AI Application Intern
% Compile: pdflatex CV_Nguyen_Khac_Bao_Asoft_AI_Automation_Intern.tex
%------------------------------------------------------

\documentclass[10pt,a4paper]{article}

\usepackage[T1]{fontenc}
\usepackage[utf8]{inputenc}
\usepackage[a4paper,top=0.32in,bottom=0.32in,left=0.46in,right=0.46in]{geometry}
\usepackage{enumitem}
\usepackage{hyperref}
\usepackage{titlesec}
\usepackage{microtype}
\usepackage{xcolor}

\hypersetup{
    colorlinks=true,
    urlcolor=blue,
    linkcolor=black,
    pdfborder={0 0 0}
}

\titleformat{\section}
    {\normalfont\large\bfseries}
    {}
    {0em}
    {}
    [\vspace{1pt}\titlerule]

\titlespacing*{\section}{0pt}{2pt}{1pt}

\setlength{\parindent}{0pt}
\setlength{\parskip}{0pt}
\pagestyle{empty}
\sloppy

\newenvironment{cvitems}{%
    \begin{itemize}[leftmargin=1.1em,itemsep=0pt,topsep=0pt,parsep=0pt]
}{%
    \end{itemize}
}

\newcommand{\entryhead}[2]{%
    \noindent\textbf{#1}\hfill #2\\[-2pt]
}

\newcommand{\entrysub}[1]{%
    \noindent\textit{#1}\\[-1pt]
}

\begin{document}
\fontsize{9}{9.9}\selectfont

\begin{center}
    {\Large\textbf{NGUYEN KHAC BAO}}\\[-1pt]
    AI Engineer Intern -- LLM/RAG \& ERP Workflow Integration\\[2pt]
    Ho Chi Minh City, Vietnam
    \,\textbar\,
    \href{mailto:nguyenkhacbaowork564@gmail.com}{nguyenkhacbaowork564@gmail.com}
    \,\textbar\,
    (+84) 039 586 4429
    \,\textbar\,
    \href{https://github.com/NguyenKhacBao564}{github.com/NguyenKhacBao564}
\end{center}

\section{CAREER OBJECTIVE}
AI Engineer Intern candidate focused on building practical AI tools for business operations, ERP-style workflows, document processing, reporting, and chatbot automation. Hands-on with Python, FastAPI, REST APIs, Streamlit, Gemini API, RAG chatbot, SQL/vector databases, Docker, and cloud deployment.

\section{EDUCATION}
\entryhead{Posts and Telecommunications Institute of Technology (PTIT)}{2022 -- 2027}
\entrysub{Bachelor of Information Technology}
Relevant Coursework: Machine Learning, Deep Learning, Computer Vision.

\section{TECHNICAL SKILLS}
\textbf{Programming:} Python, SQL basics, JavaScript basics\\
\textbf{LLM / RAG:} Gemini API, OpenAI-compatible APIs, RAG pipelines, embeddings, vector search, prompt engineering\\
\textbf{Backend / API:} FastAPI, REST API, Streamlit, JSON/API integration\\
\textbf{Data / SQL:} Pandas, CSV/JSONL processing, MySQL/Cloud SQL, MariaDB, Qdrant Cloud\\
\textbf{DevOps / Cloud:} Docker, Linux, Git/GitHub, Google Cloud Run, AWS S3/SageMaker, Vast.ai, Hugging Face\\
\textbf{AI / CV / OCR:} OpenCV, OCR, YOLOv8n, object detection, dataset QA

\section{LANGUAGES}
Vietnamese: Native \textbar\ English: TOEIC 610/990; technical reading proficiency for documentation, API references, model cards, and engineering articles.

\section{PROJECTS}
\entryhead{Vietnamese Legal RAG Chatbot}{2026}
\entrysub{Personal Project \textbar\ Python, FastAPI, Streamlit, Gemini API, Qdrant Cloud, Cloud SQL, Docker, Google Cloud Run \textbar\ Source: \href{https://github.com/NguyenKhacBao564/legal_RaG_Chatbot_VietNamese-System}{GitHub}}
\begin{cvitems}
    \item Built and deployed a \textbf{RAG chatbot} with \textbf{Streamlit} frontend and \textbf{FastAPI} backend as separate \textbf{Cloud Run} services.
    \item Indexed 32,980 Vietnamese QA/context records into \textbf{Qdrant Cloud} using \textbf{Gemini embeddings} for retrieval-grounded answers.
    \item Implemented \textbf{chat APIs}, \textbf{Cloud SQL} chat history, and Secret Manager-based credential handling.
    \item Designed the architecture to support internal documents, SQL-backed records, ERP-style workflows, async indexing, and optional OpenAI-compatible LLM serving.
\end{cvitems}

\entryhead{AutoFall Labeler Tool}{2026}
\entrysub{Personal Project \textbar\ Python, Streamlit, OpenCV, YOLO-Pose, Pandas \textbar\ Source: \href{https://github.com/NguyenKhacBao564/autofall-labeler-tool}{GitHub}}
\begin{cvitems}
    \item Built a \textbf{human-in-the-loop workflow tool} with Streamlit review UI, metadata/annotation management, YOLO-Pose keypoint suggestions, and reviewed-only export.
    \item Added evaluation reports for \textbf{auto-label accuracy}, correction rate, pose missing rate, and keypoint edit rate.
    \item Designed CSV-based source-of-truth records for annotations, metadata, keypoints, and review events to support reproducible workflow documentation.
\end{cvitems}

\entryhead{Mini Microservices Social Web}{2026}
\entrysub{Personal Project \textbar\ React, Node.js, Express Gateway, RabbitMQ, PostgreSQL/Cloud SQL, Docker Compose \textbar\ Source: \href{https://github.com/NguyenKhacBao564/Mini_Microservices_Socail_Web}{GitHub}}
\begin{cvitems}
    \item Built a mini social web system with \textbf{API Gateway}, user/post/comment/feed/notification services, JWT-protected routes, and \textbf{RabbitMQ} event messaging.
    \item Containerized frontend, gateway, backend services, RabbitMQ, and database configuration using \textbf{Docker Compose}.
\end{cvitems}

\entryhead{Traffic Violation Detection \& License Plate Recognition}{2026}
\entrysub{Personal Project \textbar\ Python, OpenCV, YOLO, OCR, JSON \textbar\ Source: \href{https://github.com/NguyenKhacBao564/Traffic-Violation-Detection-and-License-Plate-Recognition}{GitHub}}
\begin{cvitems}
    \item Built an AI video analytics pipeline for vehicle detection/tracking, traffic-light state estimation, stop-line violation logic, \textbf{OCR}, and JSON evidence generation.
    \item Fine-tuned a YOLO plate detector on 8,598 CCPD images, achieving 0.989 precision, 0.998 recall, and 0.994 mAP50.
\end{cvitems}

\section{RELEVANT STRENGTHS}
\begin{cvitems}
    \item Built practical AI automation tools using Python, Streamlit/FastAPI, API integration, and Docker-based deployment.
    \item Comfortable processing CSV/JSON data, designing metadata schemas, generating QA reports, and documenting workflow outcomes.
    \item Demonstrated problem-solving through real-world CV/OCR pipelines involving tracking, noisy data, edge cases, and human-in-the-loop validation.
\end{cvitems}

\section{ACTIVITIES}
\entryhead{AI Research Group -- University Club}{2025 -- Present}
\entrysub{Member}
\begin{cvitems}
    \item Discussed AI applications, workflow automation ideas, Machine Learning, Deep Learning, Computer Vision, and RAG/LLM systems.
\end{cvitems}

\end{document}
